\documentclass[12pt]{article}
\usepackage[utf8]{inputenc}
\usepackage{amsmath, amssymb}
\usepackage{geometry}
\geometry{margin=1in}

\title{The Specific-Factors Model}
\author{ECON 308 -- International Economic Relations}
\date{Fall 2025}

\begin{document}
\maketitle

\section*{Problem Statement}

Consider two countries, \textbf{Home} and \textbf{Foreign}, each producing two goods: \textbf{Manufactures (M)} and \textbf{Food (F)}. 

\begin{itemize}
\item Labor ($L$) is mobile between sectors within each country but not internationally.
\item Capital ($K$) is specific to manufacturing, and land ($T$) is specific to food.
\item Production functions are:
\[
Q_M = K^{1/2} L_M^{1/2}, 
\quad 
Q_F = T^{1/2} L_F^{1/2}, 
\quad 
L_M + L_F = L.
\]
\end{itemize}

Endowments:
\[
\text{Home: } K=144,\; T=64,\; L=100 
\quad \quad 
\text{Foreign: } K^{*}=64,\; T^{*}=144,\; L^{*}=100.
\]

Assume competitive markets.

\bigskip

\section*{Part A. Autarky Equilibrium}

\subsection*{Step 1: Marginal Products of Labor}
Differentiate production functions:
\[
MPL_M = \frac{\partial Q_M}{\partial L_M} = \tfrac12 K^{1/2} L_M^{-1/2}, 
\quad 
MPL_F = \tfrac12 T^{1/2} L_F^{-1/2}.
\]
\textit{Interpretation:} Both marginal products decline as labor increases in the sector (diminishing marginal returns). 

\subsection*{Step 2: Wage Equalization Condition}
Since labor is mobile, wages must be equal across sectors:
\[
P_M \cdot MPL_M = P_F \cdot MPL_F.
\]
This determines the division of labor between $M$ and $F$.

\subsection*{Step 3: Home in Autarky}
Let $P_M=2$, $P_F=1$. Home has $K^{1/2}=12$, $T^{1/2}=8$.  
Condition:
\[
2 \cdot \frac{12}{\sqrt{L_M}} = \frac{8}{\sqrt{100-L_M}}.
\]
Simplify:
\[
\frac{24}{\sqrt{L_M}} = \frac{8}{\sqrt{100-L_M}}.
\]
Square both sides:
\[
\frac{576}{L_M} = \frac{64}{100-L_M}.
\]
Cross-multiply:
\[
576(100-L_M) = 64 L_M \quad \Rightarrow \quad 57,600 - 576L_M = 64L_M.
\]
\[
57,600 = 640L_M \quad \Rightarrow \quad L_M = 90,\; L_F = 10.
\]
The wage:
\[
w = P_M \cdot MPL_M = 2 \cdot \frac{12}{2\sqrt{90}} = \frac{12}{\sqrt{90}} \approx 1.265.
\]
\textit{Explanation:} Most labor is allocated to manufacturing since Home is capital-abundant.

\subsection*{Step 4: Foreign in Autarky}
Foreign has $K^{*1/2}=8$, $T^{*1/2}=12$.  
Condition:
\[
2 \cdot \frac{8}{\sqrt{L_M^{*}}} = \frac{12}{\sqrt{100-L_M^{*}}}.
\]
Simplify:
\[
\frac{16}{\sqrt{L_M^{*}}} = \frac{12}{\sqrt{100-L_M^{*}}}.
\]
Square both sides:
\[
\frac{256}{L_M^{*}} = \frac{144}{100-L_M^{*}}.
\]
Cross-multiply:
\[
256(100-L_M^{*}) = 144 L_M^{*}.
\]
\[
25,600 - 256 L_M^{*} = 144 L_M^{*} \quad \Rightarrow \quad 25,600 = 400 L_M^{*}.
\]
\[
L_M^{*} = 64,\; L_F^{*} = 36.
\]
The wage:
\[
w^{*} = P_M \cdot MPL_M^{*} = 2 \cdot \frac{8}{2\sqrt{64}} = \frac{8}{8} = 1.
\]
\textit{Explanation:} Foreign allocates relatively more labor to food, since it is land-abundant.

\subsection*{Step 5: Autarky Relative Prices}
- In Home, manufactures are relatively cheap (low $P_M/P_F$).  
- In Foreign, food is relatively cheap.  
\textit{Implication:} Home will export $M$, Foreign will export $F$ once trade is allowed.

\bigskip

\section*{Part B. Opportunity Costs and Comparative Advantage}

\subsection*{Step 6: Opportunity Cost Formula}
A reallocation of one unit of labor from $F$ to $M$ increases $Q_M$ by $MPL_M$ and reduces $Q_F$ by $MPL_F$.  
Hence:
\[
OC_M = \frac{MPL_F}{MPL_M}.
\]

\subsection*{Step 7: Numerical Calculation at $L_M=L_F=50$}
- Home: $OC_M = \tfrac{8}{12} = \tfrac{2}{3} \approx 0.67$.  
- Foreign: $OC_M^{*} = \tfrac{12}{8} = 1.5$.  

\subsection*{Step 8: Comparative Advantage}
- Home’s $OC_M$ is lower $\Rightarrow$ comparative advantage in manufactures.  
- Foreign’s $OC_M^{*}$ is higher $\Rightarrow$ comparative advantage in food.  

\textit{Conclusion:} The pattern of trade is consistent with relative factor abundance.

\bigskip

\section*{Part C. Free Trade Equilibrium}

\subsection*{Step 9: Home with World Price $P_M/P_F=1.2$}
Condition:
\[
1.2 \cdot \frac{12}{\sqrt{L_M}} = \frac{8}{\sqrt{100-L_M}}.
\]
Simplify:
\[
\frac{14.4}{\sqrt{L_M}} = \frac{8}{\sqrt{100-L_M}}.
\]
Square both sides:
\[
\frac{207.36}{L_M} = \frac{64}{100-L_M}.
\]
\[
207.36(100-L_M) = 64 L_M.
\]
\[
20,736 - 207.36 L_M = 64 L_M \quad \Rightarrow \quad 20,736 = 271.36 L_M.
\]
\[
L_M \approx 76.4, \; L_F \approx 23.6.
\]
Wage:
\[
w = P_M \cdot MPL_M = 1.2 \cdot \frac{12}{2\sqrt{76.4}} \approx 0.825.
\]
\textit{Interpretation:} Compared to autarky, more labor shifts into manufacturing, but not as extreme as before because the world price is only moderately above 1.

\subsection*{Step 10: Foreign with World Price $1.2$}
Condition:
\[
1.2 \cdot \frac{8}{\sqrt{L_M^{*}}} = \frac{12}{\sqrt{100-L_M^{*}}}.
\]
Simplify:
\[
\frac{9.6}{\sqrt{L_M^{*}}} = \frac{12}{\sqrt{100-L_M^{*}}}.
\]
Square both sides:
\[
\frac{92.16}{L_M^{*}} = \frac{144}{100-L_M^{*}}.
\]
\[
92.16(100-L_M^{*}) = 144 L_M^{*}.
\]
\[
9,216 - 92.16 L_M^{*} = 144 L_M^{*} \quad \Rightarrow \quad 9,216 = 236.16 L_M^{*}.
\]
\[
L_M^{*} \approx 39.0,\; L_F^{*} \approx 61.0.
\]
Wage:
\[
w^{*} = 1.2 \cdot \frac{8}{2\sqrt{39}} \approx 0.769.
\]

\subsection*{Step 11: Outputs}
- Home: 
\[
Q_M = 12\sqrt{76.4} \approx 104.7, 
\quad 
Q_F = 8\sqrt{23.6} \approx 39.0.
\]
- Foreign: 
\[
Q_M^{*} = 8\sqrt{39.0} \approx 49.9, 
\quad 
Q_F^{*} = 12\sqrt{61.0} \approx 93.6.
\]

\subsection*{Step 12: Trade Pattern}
- Home produces relatively more $M$, exports $M$, imports $F$.  
- Foreign produces relatively more $F$, exports $F$, imports $M$.  

\bigskip

\section*{Part D. Distributional Effects of Trade}

\subsection*{Step 13: Real Wage in Home}
- In terms of $M$: $\tfrac{w}{P_M} = \tfrac{0.825}{1.2} \approx 0.69$ (lower than in autarky, workers lose relative to $M$).  
- In terms of $F$: $\tfrac{w}{P_F} = 0.825$ (higher than before, workers gain relative to $F$).  

\textit{Workers face ambiguous welfare changes.}

\subsection*{Step 14: Returns to Capital and Land in Home}
- Capital: more labor in $M$ raises its marginal product, capital owners gain.  
- Land: less labor in $F$ lowers land productivity, landowners lose.  

\subsection*{Step 15: Distribution in Foreign}
- Workers: similar ambiguity as in Home.  
- Capital (specific to $M$): loses, since less labor allocated to manufacturing.  
- Land (specific to $F$): gains, as labor shifts into food.  

\bigskip

\section*{Summary}
\begin{itemize}
\item Home is capital-abundant $\Rightarrow$ comparative advantage in manufactures.  
\item Foreign is land-abundant $\Rightarrow$ comparative advantage in food.  
\item Trade benefits the specific factor used intensively in the export sector (capital in Home, land in Foreign).  
\item Trade hurts the specific factor used in the import sector (land in Home, capital in Foreign).  
\item Labor is mobile but experiences ambiguous welfare effects, depending on consumption patterns.  
\end{itemize}

\end{document}
